\documentclass[a4paper,12pt]{article}
\usepackage[utf8]{inputenc}
\usepackage[polish]{babel}
\usepackage[T1]{fontenc}
\usepackage{amsmath}
\usepackage{graphicx}
\usepackage{geometry}
\usepackage{float}
\usepackage{caption}
\usepackage{listings}
\usepackage{xcolor}
\usepackage{subfigure}
\usepackage{booktabs}
\usepackage{array}

\geometry{margin=2.5cm}

% Styl dla kodu MATLAB
\lstdefinestyle{matlab}{
    language=Matlab,
    basicstyle=\ttfamily\small,
    keywordstyle=\color{blue}\bfseries,
    commentstyle=\color{gray}\itshape,
    stringstyle=\color{orange},
    numbers=left,
    numberstyle=\tiny\color{gray},
    numbersep=8pt,
    frame=single,
    breaklines=true,
    backgroundcolor=\color{gray!8},
    captionpos=b
}

\title{
\Large Zaawansowane metody przetwarzania i analizy sygnałów \\[6pt]
\textbf{Raport z postępów – Analiza HRV i klasyfikacja grup wiekowych}
}

\author{Piotr Sawicki 319003, Michał Odziemkowski 311267}
\date{Maj 2026}

\begin{document}

\maketitle

% =========================================================

\section{Cel projektu}
Głównym celem niniejszego projektu jest przeprowadzenie zaawansowanej analizy częstotliwościowej sygnału zmienności rytmu zatokowego (ang. \textit{Heart Rate Variability} – HRV)~. Projekt zakłada identyfikację istotnych statystycznie różnic w parametrach widmowych pomiędzy osobami młodymi a seniorami. Efektem końcowym ma być stworzenie stabilnego modelu klasyfikacyjnego, zdolnego do automatycznego rozróżniania grup wiekowych na podstawie cech wyekstrahowanych z sygnału EKG.


\section{Zakres wykonanych prac}

W ramach bieżącego etapu projektu zrealizowano następujące zadania:

\begin{itemize}
    \item Instalacja pakietu \textit{WFDB Toolbox for MATLAB} umożliwiającego wczytywanie danych z bazy PhysioNet.
    \item Zapoznanie się ze strukturą bazy \textit{Fantasia Database} oraz formatem przechowywanych plików.
    \item Wczytanie danych z bazy przy użyciu funkcji \texttt{rdsamp} oraz \texttt{rdann}.
    \item Ekstrakcja sygnału EKG oraz detekcja załamków R.
    \item Wyznaczenie odstępów RR oraz ich reprezentacji czasowej.
    \item Implementacja i porównanie dwóch metod estymacji widma mocy: FFT oraz Welch.
    \item Wyznaczenie parametrów częstotliwościowych: $LF$, $HF$ oraz $LF/HF$.
    \item Przygotowanie zbioru cech do klasyfikacji.
\end{itemize}

% =========================================================

\section{Środowisko i narzędzia}

\subsection{WFDB Toolbox for MATLAB}

Dane z bazy \textit{Fantasia Database} przechowywane są w formacie WFDB (\textit{WaveForm DataBase}) -- własnym formacie binarnym platformy PhysioNet, który nie jest natywnie obsługiwany przez środowisko MATLAB. W celu umożliwienia odczytu tych danych zainstalowano pakiet \textbf{WFDB Toolbox for MATLAB}, dostępny bezpłatnie na stronie \texttt{physionet.org}.

Pakiet dostarcza dwie kluczowe funkcje wykorzystywane w projekcie:

\begin{itemize}
    \item \texttt{rdsamp} - wczytuje próbki sygnału fizjologicznego (EKG) z pliku binarnego \texttt{.dat} wraz z częstotliwością próbkowania odczytaną z nagłówka \texttt{.hea}:

    \item \texttt{rdann} - wczytuje adnotacje ekspertów z pliku \texttt{.ecg}, zawierające numery próbek odpowiadające pozycjom załamków R:

\end{itemize}

% =========================================================

\section{Baza danych Fantasia Database}

\subsection{Opis bazy}

Baza \textit{Fantasia Database} udostępniona jest na platformie PhysioNet i zawiera długoterminowe nagrania sygnałów fizjologicznych wykonane u zdrowych ochotników w dwóch grupach wiekowych:

\begin{itemize}
    \item \textbf{Young} -- 20 młodych osób (wiek 21--34 lat),
    \item \textbf{Elderly} -- 20 starszych osób (wiek 68--85 lat).
\end{itemize}

Każde nagranie trwa \textbf{120 minut} i zostało zarejestrowane podczas oglądania przez badanego wybranego filmu fantasy. Baza zawiera łącznie \textbf{40 rekordów}.

\subsection{Struktura plików}

Każdy rekord w bazie składa się z trzech plików o różnych rozszerzeniach, pełniących odmienne funkcje:

\begin{table}[H]
\centering
\caption{Rodzaje plików wchodzących w skład bazy Fantasia Database}
\begin{tabular}{@{}cll@{}}
\toprule
\textbf{Rozszerzenie} & \textbf{Zawartość} & \textbf{Wykorzystanie w projekcie} \\
\midrule
\texttt{.dat} & Próbki sygnału EKG (format binarny) & Wczytanie przez \texttt{rdsamp} \\
\texttt{.hea} & Nagłówek: $f_s$, liczba kanałów, czas trwania & Automatycznie przez \texttt{rdsamp} \\
\texttt{.ecg} & Adnotacje: pozycje załamków R & Wczytanie przez \texttt{rdann} \\
\bottomrule
\end{tabular}
\end{table}

\subsection{Konwencja nazw rekordów}

Nazwy rekordów zakodowane są w następujący sposób:

\begin{center}
\texttt{f\{seria\}\{grupa\}\{numer\}}
\end{center}

gdzie seria przyjmuje wartości \texttt{1} lub \texttt{2}, grupa to \texttt{y} (young) lub \texttt{o} (old/elderly), a numer to dwucyfrowy identyfikator pacjenta od \texttt{01} do \texttt{10}. Przykładowo rekord \texttt{f1y01} oznacza serię~1, grupę young, pacjent~1, natomiast \texttt{f2o05} to seria~2, grupa elderly, pacjent~5.

% =========================================================

\section{Implementacja}

Analiza została przeprowadzona w środowisku MATLAB. Proces obliczeniowy dla każdego rekordu przebiega według następujących kroków:

\begin{enumerate}

    \item \textbf{Wyznaczenie odstępów RR.}
    Odstępy RR obliczono jako różnice czasów kolejnych załamków R:
    \[
        RR_i = t_{i+1} - t_i \quad [s]
    \]

    \item \textbf{Interpolacja sygnału RR.}
    Sygnał RR jest nierównomiernie próbkowany (każdy odstęp ma inną długość), dlatego przed analizą widmową dokonano interpolacji metodą \textit{spline} do równomiernej siatki czasowej o częstotliwości $f_s = 4$~Hz. Wartość ta zapewnia spełnienie twierdzenia Nyquista dla składowej HF ($f_{max} = 0{,}4$~Hz):
    \[
        f_s = 4 \text{ Hz} > 2 \cdot f_{max} = 0{,}8 \text{ Hz}
    \]

    \item \textbf{Usunięcie trendu.}
    Przed estymacją widma zastosowano funkcję \texttt{detrend} w celu usunięcia wolnozmiennego trendu (składowej stałej i liniowej), który mógłby zaburzać wyniki w paśmie LF.

    \item \textbf{Estymacja widma mocy.}
    Gęstość widmowa mocy (PSD) wyznaczona została dwiema metodami -- FFT oraz Welch -- których porównanie opisano w sekcji~\ref{sec:porownanie}.

    \item \textbf{Wyznaczenie mocy w pasmach.}
    Moc w pasmach LF i HF obliczono metodą trapezów (\texttt{trapz}):
    \[
        LF = \int_{0{,}04}^{0{,}15} PSD(f)\,df, \qquad
        HF = \int_{0{,}15}^{0{,}40} PSD(f)\,df
    \]

    \item \textbf{Obliczenie wskaźnika $LF/HF$.}
    \[
        \frac{LF}{HF}
    \]
    Wskaźnik ten odzwierciedla równowagę pomiędzy układem współczulnym (LF) a przywspółczulnym (HF).

\end{enumerate}

% =========================================================

\section{Porównanie metod estymacji PSD: FFT vs Welch}
\label{sec:porownanie}

\subsection{Opis metod}

\textbf{FFT (periodogram)} wyznacza widmo z całego sygnału jednocześnie przy użyciu okna Hanna (w celu redukcji przecieku widmowego). Metoda charakteryzuje się wysoką rozdzielczością częstotliwościową, lecz wysoką wariancją estymatora -- widmo jest nieregularne i wrażliwe na artefakty.

\textbf{Metoda Welcha} dzieli sygnał na nakładające się segmenty, oblicza widmo każdego z nich, a następnie uśrednia wyniki. Dzięki temu wariancja estymatora jest istotnie mniejsza kosztem nieznacznego obniżenia rozdzielczości częstotliwościowej.

\subsection{Implementacja w MATLAB}

\begin{lstlisting}[style=matlab, caption={Estymacja PSD metodą FFT}]
N   = length(RR_interp);
win = hann(N)';
X   = fft(RR_interp .* win);
PSD = (1/(fs*N)) * abs(X).^2;
PSD = PSD(1:floor(N/2)+1);
PSD(2:end-1) = 2 * PSD(2:end-1);   % widmo jednostronne
f   = (0:floor(N/2)) * fs / N;
\end{lstlisting}

\begin{lstlisting}[style=matlab, caption={Estymacja PSD metodą Welcha}]
[PSD, f] = pwelch(RR_interp, [], [], [], fs);
\end{lstlisting}

\subsection{Wyniki porównania}

\begin{figure}[H]
\centering
\includegraphics[width=\textwidth]{../wykresy/wykres_ekg_rpeaks.pdf}
\caption{Fragment sygnału EKG (pierwsze 30~s) z zaznaczonymi załamkami R dla rekordu \texttt{f1y01} (young) oraz \texttt{f1o01} (elderly).}
\end{figure}

\begin{figure}[H]
\centering
\includegraphics[width=\textwidth]{../wykresy/wykres_widmo_young.pdf}
\caption{Porównanie widm HRV wyznaczonych metodą FFT i Welcha dla rekordu \texttt{f1y01} (young). Kolorem niebieskim zaznaczono pasmo LF (0{,}04--0{,}15~Hz), czerwonym pasmo HF (0{,}15--0{,}40~Hz).}
\end{figure}

\begin{figure}[H]
\centering
\includegraphics[width=\textwidth]{../wykresy/wykres_widmo_elderly.pdf}
\caption{Porównanie widm HRV wyznaczonych metodą FFT i Welcha dla rekordu \texttt{f1o01} (elderly).}
\end{figure}

\subsection{Wybór metody}

Na podstawie przeprowadzonego porównania do dalszej analizy wybrano \textbf{metodę Welcha}. Uzasadnienie wyboru:

\begin{itemize}
    \item Nagrania trwają 120 minut -- sygnał RR jest silnie niestacjonarny (zmieniają się fazy czuwania i snu), dlatego stabilność estymatora jest ważniejsza niż rozdzielczość.
    \item Metoda Welcha daje gładkie, powtarzalne widmo, co przekłada się na stabilniejsze wartości parametrów LF, HF i $LF/HF$ -- kluczowe dla klasyfikatora SVM.
    \item Periodogram FFT jest wrażliwy na artefakty ruchowe i ektopowe uderzenia serca, które mogą wystąpić podczas tak długich nagrań.
    \item Metoda Welcha jest standardem w analizie HRV rekomendowanym przez \textit{Task Force} (European Society of Cardiology, 1996).
\end{itemize}

% =========================================================

\section{Wyniki wstępne}

Uzyskane wyniki wskazują na wyraźne różnice pomiędzy grupami:

\begin{itemize}
    \item Wyraźnie większą składową wysokoczęstotliwościową ($HF$) w grupie młodych, co świadczy o silniejszej aktywności układu przywspółczulnego.
    \item Dominację składowych niskoczęstotliwościowych ($LF$) u osób starszych.
    \item Wyższy wskaźnik $LF/HF$ w grupie elderly, zgodny z literaturą medyczną opisującą osłabienie regulacji przywspółczulnej wraz z wiekiem.
\end{itemize}

% =========================================================

\section{Klasyfikacja}

Zaimplementowano klasyfikator SVM z jądrem RBF (\textit{Radial Basis Function}).

\begin{itemize}
    \item Podział danych: 80\% trening, 20\% test
    \item Standaryzacja cech metodą Z-score (parametry wyznaczone na zbiorze treningowym)
    \item Cechy wejściowe: $[LF,\ HF,\ LF/HF]$
\end{itemize}

Wstępna skuteczność klasyfikacji wynosi:

\[
    Accuracy \approx XX\%
\]

% =========================================================

\section{Problemy napotkane podczas realizacji}

W trakcie prac napotkano następujące trudności:

\begin{itemize}
    \item Błędna konwencja nazw plików elderly w bazie -- rekordy oznaczone są literą \texttt{o} (od \textit{old}), a nie \texttt{e} (od \textit{elderly}), co początkowo uniemożliwiało ich wczytanie.
    \item Konieczność odpowiedniej normalizacji danych przed klasyfikacją -- bez standaryzacji SVM z jądrem RBF dawał istotnie gorsze wyniki.
    \item Wrażliwość wyników FFT na artefakty w długich nagraniach -- potwierdziła potrzebę wyboru metody Welcha.
\end{itemize}

% =========================================================

\section{Plan dalszych prac}

W kolejnych etapach planuje się:

\begin{itemize}
    \item Wprowadzenie analizy segmentowej (segmenty 5-minutowe) zamiast całego nagrania -- zgodnie ze standardem \textit{Task Force} dla krótkoterminowej analizy HRV.
    \item Rozszerzenie zbioru cech o parametry dziedziny czasu: SDNN, RMSSD.
    \item Optymalizację hiperparametrów klasyfikatora (grid search po $C$ i $\gamma$).
    \item Zastąpienie podziału hold-out walidacją \textit{Repeated 10-fold cross-validation} z testem statystycznym (paired t-test) w celu wiarygodnego porównania metod FFT i Welch.
    \item Analizę statystyczną różnic między grupami (test Manna-Whitneya).
\end{itemize}

% =========================================================

\section{Wnioski}

Dotychczasowe wyniki potwierdzają, że analiza widmowa HRV stanowi skuteczne narzędzie do rozróżniania grup wiekowych. Obserwowany spadek składowej $HF$ u osób starszych jest zgodny z literaturą i wynika z osłabienia układu przywspółczulnego wraz z wiekiem. Spośród dwóch testowanych metod estymacji widma -- FFT i Welch -- wybrano metodę Welcha jako stabilniejszą i bardziej odpowiednią dla długich, niestacjonarnych sygnałów biologicznych.

\end{document}